
\documentclass{article}
\usepackage{colortbl}
\usepackage{makecell}
\usepackage{multirow}
\usepackage{supertabular}

\begin{document}

\newcounter{utterance}

\centering \large Interaction Transcript for game `clean\_up', experiment `2\_hard\_3obj\_ro', episode 3 with nemotron{-}3{-}ultra.
\vspace{24pt}

{ \footnotesize  \setcounter{utterance}{1}
\setlength{\tabcolsep}{0pt}
\begin{supertabular}{c@{$\;$}|p{.15\linewidth}@{}p{.15\linewidth}p{.15\linewidth}p{.15\linewidth}p{.15\linewidth}p{.15\linewidth}}
    \# & $\;$A & \multicolumn{4}{c}{Game Master} & $\;\:$B\\
    \hline

    \theutterance \stepcounter{utterance}  
    & & \multicolumn{4}{p{0.6\linewidth}}{
        \cellcolor[rgb]{0.9,0.9,0.9}{
            \makecell[{{p{\linewidth}}}]{
                \texttt{\tiny{[P1$\langle$GM]}}
                \texttt{Eu sunt maestrul vostru de joc, iar voi jucați un joc colaborativ cu următoarea grilă ca tablă de joc:} \\
\\ 
\texttt{```} \\
\texttt{ 1234567} \\
\texttt{╔═╤═╤═══╗} \\
\texttt{║L│◌│◌◌◌║ 1} \\
\texttt{║C├─┤◌◌◌║ 2} \\
\texttt{╟─┤◌├─┬─╢ 3} \\
\texttt{║◌│◌│◌│◌║ 4} \\
\texttt{╟─┼─┘◌│◌║ 5} \\
\texttt{║◌│◌◌┌┴─╢ 6} \\
\texttt{║◌│◌◌│◌P║ 7} \\
\texttt{╚═╧══╧══╝} \\
\\ 
\texttt{```} \\
\\ 
\texttt{* Marginea superioară afișează coordonatele x crescând spre dreapta, iar marginea dreaptă coordonatele y crescând în jos.} \\
\texttt{* Următoarele obiecte sunt plasate aleatoriu pe grila voastră: 'C', 'L', 'P'.} \\
\\ 
\\ 
\texttt{Celălalt jucător vede o variantă a tablei de joc, în care obiectele sunt plasate în diferite locații aleatorii. Nu puteți vedea tabla de joc a celuilalt jucător, iar el nu o poate vedea pe a dumneavoastră.} \\
\\ 
\\ 
\texttt{**Obiectiv:**} \\
\\ 
\texttt{Ambii jucători trebuie să mute obiectele de pe fundalul lor respectiv astfel încât obiectele identice să ajungă la aceleași coordonate. Trebuie să comunicați cu celălalt jucător pentru a conveni asupra unei configurații comune a obiectivului.} \\
\\ 
\\ 
\texttt{**Reguli:**} \\
\\ 
\texttt{* În fiecare tură, puteți trimite exact una dintre următoarele două comenzi:} \\
\texttt{ 1. `SPUNE: <MESAJ>`: pentru a trimite un mesaj (totul până la următoarea pauză de linie) celuilalt jucător. Eu îl voi transmite partenerului tău.} \\
\texttt{2. `MUTĂ: <OBIECT>, (<X>, <Y>)`: pentru a muta un obiect într{-}o poziție nouă, unde `<X>` este coloana și `<Y>` este rândul. Te voi informa dacă mutarea ta a avut succes sau nu.} \\
\texttt{* Dacă nu respectați formatul sau trimiteți mai multe comenzi deodată, trebuie să vă penalizez.} \\
\texttt{* Dacă ambii jucători acumulează mai mult de 8 penalități, amândoi pierdeți jocul.} \\
\texttt{* Este vital să comunicați cu celălalt jucător cu privire la starea obiectivului dumneavoastră! *Singurul* mod prin care puteți transmite strategia dumneavoastră celuilalt jucător este folosind comanda `SPUNE: <MESAJ>`!} \\
\\ 
\\ 
\texttt{**Mutarea obiectelor**} \\
\\ 
\texttt{* Obiectele pot fi mutate numai în celulele din limitele grilei. Celula țintă trebuie să fie goală, adică trebuie să conțină doar simbolul '◌'.} \\
\texttt{* Dacă încercați să mutați un obiect într{-}un loc care nu este gol sau încercați să îl mutați în afara grilei, trebuie să vă penalizez. Aveți dreptul la o nouă încercare.} \\
\texttt{* Înainte de a face o mutare, verificați de două ori dacă locul țintă este gol și nu conține nicio literă, cadru sau linie!} \\
\\ 
\\ 
\texttt{**Sfârșitul jocului**} \\
\\ 
\texttt{Dacă credeți că ați atins obiectivul de aliniere a tuturor obiectelor, puteți cere celuilalt jucător să încheie jocul trimițând `SPUNE: terminat?`. Dacă celălalt jucător vă cere să încheiați jocul, iar dvs. răspundeți cu `SPUNE: terminat!`, jocul se va încheia.} \\
\\ 
\texttt{Ambii jucători câștigă dacă jocul se încheie în decursul a 12 runde, unde o rundă este definită ca fiind trimiterea de către doi jucători a câte unei comenzi valide.} \\
\\ 
\texttt{**Punctaj:**} \\
\\ 
\texttt{Cu cât obiectele identice sunt mai apropiate pe ambele table de joc, cu atât primești mai multe puncte. Penalizările îți reduc punctele. Poți bate recordul?} \\
\\ 
\\ 
\texttt{Vă rugăm să trimiteți un mesaj celuilalt jucător pentru a începe jocul!} \\
            }
        }
    }
    & \\ \\

    \theutterance \stepcounter{utterance}  
    & \multicolumn{4}{p{0.6\linewidth}}{
        \cellcolor[rgb]{0.9,0.9,0.9}{
            \makecell[{{p{\linewidth}}}]{
                \texttt{\tiny{[P1$\rangle$GM]}}
                \texttt{SPUNE: Salut! Sunt la (1,1) pentru L, (1,2) pentru C și (7,7) pentru P. Tu unde ai obiectele?} \\
            }
        }
    }
    & & \\ \\

    \theutterance \stepcounter{utterance}  
    & & & \multicolumn{2}{p{0.3\linewidth}}{
        \cellcolor[rgb]{0.9,0.9,0.9}{
            \makecell[{{p{\linewidth}}}]{
                \texttt{\tiny{[GM$|$GM]}}
                \texttt{SPUNE: Salut! Sunt la (1,1) pentru L, (1,2) pentru C și (7,7) pentru P. Tu unde ai obiectele?} \\
            }
        }
    }
    & & \\ \\

    \theutterance \stepcounter{utterance}  
    & & \multicolumn{4}{p{0.6\linewidth}}{
        \cellcolor[rgb]{0.9,0.9,0.9}{
            \makecell[{{p{\linewidth}}}]{
                \texttt{\tiny{[GM$\rangle$P2]}}
                \texttt{Eu sunt maestrul vostru de joc, iar voi jucați un joc colaborativ cu următoarea grilă ca tablă de joc:} \\
\\ 
\texttt{```} \\
\texttt{ 1234567} \\
\texttt{╔═╤═╤═══╗} \\
\texttt{║◌│◌│◌◌◌║ 1} \\
\texttt{║◌├─┤◌L◌║ 2} \\
\texttt{╟─┤P├─┬─╢ 3} \\
\texttt{║◌│◌│◌│◌║ 4} \\
\texttt{╟─┼─┘C│◌║ 5} \\
\texttt{║◌│◌◌┌┴─╢ 6} \\
\texttt{║◌│◌◌│◌◌║ 7} \\
\texttt{╚═╧══╧══╝} \\
\\ 
\texttt{```} \\
\\ 
\texttt{* Marginea superioară afișează coordonatele x crescând spre dreapta, iar marginea dreaptă coordonatele y crescând în jos.} \\
\texttt{* Următoarele obiecte sunt plasate aleatoriu pe grila voastră: 'C', 'L', 'P'.} \\
\\ 
\\ 
\texttt{Celălalt jucător vede o variantă a tablei de joc, în care obiectele sunt plasate în diferite locații aleatorii. Nu puteți vedea tabla de joc a celuilalt jucător, iar el nu o poate vedea pe a dumneavoastră.} \\
\\ 
\\ 
\texttt{**Obiectiv:**} \\
\\ 
\texttt{Ambii jucători trebuie să mute obiectele de pe fundalul lor respectiv astfel încât obiectele identice să ajungă la aceleași coordonate. Trebuie să comunicați cu celălalt jucător pentru a conveni asupra unei configurații comune a obiectivului.} \\
\\ 
\\ 
\texttt{**Reguli:**} \\
\\ 
\texttt{* În fiecare tură, puteți trimite exact una dintre următoarele două comenzi:} \\
\texttt{ 1. `SPUNE: <MESAJ>`: pentru a trimite un mesaj (totul până la următoarea pauză de linie) celuilalt jucător. Eu îl voi transmite partenerului tău.} \\
\texttt{2. `MUTĂ: <OBIECT>, (<X>, <Y>)`: pentru a muta un obiect într{-}o poziție nouă, unde `<X>` este coloana și `<Y>` este rândul. Te voi informa dacă mutarea ta a avut succes sau nu.} \\
\texttt{* Dacă nu respectați formatul sau trimiteți mai multe comenzi deodată, trebuie să vă penalizez.} \\
\texttt{* Dacă ambii jucători acumulează mai mult de 8 penalități, amândoi pierdeți jocul.} \\
\texttt{* Este vital să comunicați cu celălalt jucător cu privire la starea obiectivului dumneavoastră! *Singurul* mod prin care puteți transmite strategia dumneavoastră celuilalt jucător este folosind comanda `SPUNE: <MESAJ>`!} \\
\\ 
\\ 
\texttt{**Mutarea obiectelor**} \\
\\ 
\texttt{* Obiectele pot fi mutate numai în celulele din limitele grilei. Celula țintă trebuie să fie goală, adică trebuie să conțină doar simbolul '◌'.} \\
\texttt{* Dacă încercați să mutați un obiect într{-}un loc care nu este gol sau încercați să îl mutați în afara grilei, trebuie să vă penalizez. Aveți dreptul la o nouă încercare.} \\
\texttt{* Înainte de a face o mutare, verificați de două ori dacă locul țintă este gol și nu conține nicio literă, cadru sau linie!} \\
\\ 
\\ 
\texttt{**Sfârșitul jocului**} \\
\\ 
\texttt{Dacă credeți că ați atins obiectivul de aliniere a tuturor obiectelor, puteți cere celuilalt jucător să încheie jocul trimițând `SPUNE: terminat?`. Dacă celălalt jucător vă cere să încheiați jocul, iar dvs. răspundeți cu `SPUNE: terminat!`, jocul se va încheia.} \\
\\ 
\texttt{Ambii jucători câștigă dacă jocul se încheie în decursul a 12 runde, unde o rundă este definită ca fiind trimiterea de către doi jucători a câte unei comenzi valide.} \\
\\ 
\texttt{**Punctaj:**} \\
\\ 
\texttt{Cu cât obiectele identice sunt mai apropiate pe ambele table de joc, cu atât primești mai multe puncte. Penalizările îți reduc punctele. Poți bate recordul?} \\
\\ 
\texttt{Celălalt jucător a început jocul trimițând acest mesaj:} \\
\\ 
\texttt{"Salut! Sunt la (1,1) pentru L, (1,2) pentru C și (7,7) pentru P. Tu unde ai obiectele?"} \\
\\ 
\texttt{Care este prima ta comandă?} \\
            }
        }
    }
    & \\ \\

    \theutterance \stepcounter{utterance}  
    & & & \multicolumn{4}{p{0.6\linewidth}}{
        \cellcolor[rgb]{0.9,0.9,0.9}{
            \makecell[{{p{\linewidth}}}]{
                \texttt{\tiny{[GM$\langle$P2]}}
                \texttt{SPUNE: Salut! Am L la (2,7), C la (6,5) și P la (2,3). Propun să ne întâlnim pe coloana 4: L la (4,3), C la (4,5), P la (4,4). Sunteți de acord?} \\
            }
        }
    }
     \\ \\

    \theutterance \stepcounter{utterance}  
    & & & \multicolumn{2}{p{0.3\linewidth}}{
        \cellcolor[rgb]{0.9,0.9,0.9}{
            \makecell[{{p{\linewidth}}}]{
                \texttt{\tiny{[GM$|$GM]}}
                \texttt{SPUNE: Salut! Am L la (2,7), C la (6,5) și P la (2,3). Propun să ne întâlnim pe coloana 4: L la (4,3), C la (4,5), P la (4,4). Sunteți de acord?} \\
            }
        }
    }
    & & \\ \\

    \theutterance \stepcounter{utterance}  
    & & \multicolumn{4}{p{0.6\linewidth}}{
        \cellcolor[rgb]{0.9,0.9,0.9}{
            \makecell[{{p{\linewidth}}}]{
                \texttt{\tiny{[P1$\langle$GM]}}
                \texttt{Mesajul dvs. a fost transmis celuilalt jucător.} \\
\\ 
\\ 
\texttt{În prezent jucați runda 1 din maximum 12.} \\
\texttt{Ați acumulat în mod colectiv 0 din 8 penalizări.} \\
\texttt{Celălalt jucător a trimis acest mesaj:} \\
\\ 
\\ 
\texttt{"Salut! Am L la (2,7), C la (6,5) și P la (2,3). Propun să ne întâlnim pe coloana 4: L la (4,3), C la (4,5), P la (4,4). Sunteți de acord?"} \\
\\ 
\\ 
\texttt{Care este următoarea dvs. comandă?} \\
            }
        }
    }
    & \\ \\

    \theutterance \stepcounter{utterance}  
    & \multicolumn{4}{p{0.6\linewidth}}{
        \cellcolor[rgb]{0.9,0.9,0.9}{
            \makecell[{{p{\linewidth}}}]{
                \texttt{\tiny{[P1$\rangle$GM]}}
                \texttt{SPUNE: De acord cu coloana 4. Mă întorc cu confirmarea după ce mut L{-}ul.} \\
            }
        }
    }
    & & \\ \\

    \theutterance \stepcounter{utterance}  
    & & & \multicolumn{2}{p{0.3\linewidth}}{
        \cellcolor[rgb]{0.9,0.9,0.9}{
            \makecell[{{p{\linewidth}}}]{
                \texttt{\tiny{[GM$|$GM]}}
                \texttt{SPUNE: De acord cu coloana 4. Mă întorc cu confirmarea după ce mut L{-}ul.} \\
            }
        }
    }
    & & \\ \\

    \theutterance \stepcounter{utterance}  
    & & \multicolumn{4}{p{0.6\linewidth}}{
        \cellcolor[rgb]{0.9,0.9,0.9}{
            \makecell[{{p{\linewidth}}}]{
                \texttt{\tiny{[GM$\rangle$P2]}}
                \texttt{Mesajul dvs. a fost transmis celuilalt jucător.} \\
\\ 
\\ 
\texttt{În prezent jucați runda 2 din maximum 12.} \\
\texttt{Ați acumulat în mod colectiv 0 din 8 penalizări.} \\
\texttt{Celălalt jucător a trimis acest mesaj:} \\
\\ 
\\ 
\texttt{"De acord cu coloana 4. Mă întorc cu confirmarea după ce mut L{-}ul."} \\
\\ 
\\ 
\texttt{Care este următoarea dvs. comandă?} \\
            }
        }
    }
    & \\ \\

    \theutterance \stepcounter{utterance}  
    & & & \multicolumn{4}{p{0.6\linewidth}}{
        \cellcolor[rgb]{0.9,0.9,0.9}{
            \makecell[{{p{\linewidth}}}]{
                \texttt{\tiny{[GM$\langle$P2]}}
                \texttt{SPUNE: Înțeles. Mă pregătesc să mut P{-}ul primul la (4,4), pentru a lăsa (2,3) liber. Vă anunț când gata.} \\
            }
        }
    }
     \\ \\

    \theutterance \stepcounter{utterance}  
    & & & \multicolumn{2}{p{0.3\linewidth}}{
        \cellcolor[rgb]{0.9,0.9,0.9}{
            \makecell[{{p{\linewidth}}}]{
                \texttt{\tiny{[GM$|$GM]}}
                \texttt{SPUNE: Înțeles. Mă pregătesc să mut P{-}ul primul la (4,4), pentru a lăsa (2,3) liber. Vă anunț când gata.} \\
            }
        }
    }
    & & \\ \\

    \theutterance \stepcounter{utterance}  
    & & \multicolumn{4}{p{0.6\linewidth}}{
        \cellcolor[rgb]{0.9,0.9,0.9}{
            \makecell[{{p{\linewidth}}}]{
                \texttt{\tiny{[P1$\langle$GM]}}
                \texttt{Mesajul dvs. a fost transmis celuilalt jucător.} \\
\\ 
\\ 
\texttt{În prezent jucați runda 2 din maximum 12.} \\
\texttt{Ați acumulat în mod colectiv 0 din 8 penalizări.} \\
\texttt{Celălalt jucător a trimis acest mesaj:} \\
\\ 
\\ 
\texttt{"Înțeles. Mă pregătesc să mut P{-}ul primul la (4,4), pentru a lăsa (2,3) liber. Vă anunț când gata."} \\
\\ 
\\ 
\texttt{Care este următoarea dvs. comandă?} \\
            }
        }
    }
    & \\ \\

    \theutterance \stepcounter{utterance}  
    & \multicolumn{4}{p{0.6\linewidth}}{
        \cellcolor[rgb]{0.9,0.9,0.9}{
            \makecell[{{p{\linewidth}}}]{
                \texttt{\tiny{[P1$\rangle$GM]}}
                \texttt{MUTĂ: L, (4, 3)} \\
            }
        }
    }
    & & \\ \\

    \theutterance \stepcounter{utterance}  
    & & & \multicolumn{2}{p{0.3\linewidth}}{
        \cellcolor[rgb]{0.9,0.9,0.9}{
            \makecell[{{p{\linewidth}}}]{
                \texttt{\tiny{[GM$|$GM]}}
                \texttt{MUTĂ: L, (4, 3)} \\
            }
        }
    }
    & & \\ \\

    \theutterance \stepcounter{utterance}  
    & & & \multicolumn{2}{p{0.3\linewidth}}{
        \cellcolor[rgb]{0.9,0.9,0.9}{
            \makecell[{{p{\linewidth}}}]{
                \texttt{\tiny{[GM$|$GM]}}
                \texttt{Penalizare: (4,3) nu este gol, ci conține '├'.} \\
\\ 
\texttt{Ați acumulat în mod colectiv 1 din 8 penalizări.} \\
\texttt{Vă rugăm să încercați din nou!} \\
            }
        }
    }
    & & \\ \\

    \theutterance \stepcounter{utterance}  
    & & \multicolumn{4}{p{0.6\linewidth}}{
        \cellcolor[rgb]{0.9,0.9,0.9}{
            \makecell[{{p{\linewidth}}}]{
                \texttt{\tiny{[P1$\langle$GM]}}
                \texttt{Penalizare: (4,3) nu este gol, ci conține '├'.} \\
\\ 
\texttt{Ați acumulat în mod colectiv 1 din 8 penalizări.} \\
\texttt{Vă rugăm să încercați din nou!} \\
            }
        }
    }
    & \\ \\

    \theutterance \stepcounter{utterance}  
    & \multicolumn{4}{p{0.6\linewidth}}{
        \cellcolor[rgb]{0.9,0.9,0.9}{
            \makecell[{{p{\linewidth}}}]{
                \texttt{\tiny{[P1$\rangle$GM]}}
                \texttt{SPUNE: (4,3) e zid pe tabla mea. Încerc (4,4) pentru L.} \\
            }
        }
    }
    & & \\ \\

    \theutterance \stepcounter{utterance}  
    & & & \multicolumn{2}{p{0.3\linewidth}}{
        \cellcolor[rgb]{0.9,0.9,0.9}{
            \makecell[{{p{\linewidth}}}]{
                \texttt{\tiny{[GM$|$GM]}}
                \texttt{SPUNE: (4,3) e zid pe tabla mea. Încerc (4,4) pentru L.} \\
            }
        }
    }
    & & \\ \\

    \theutterance \stepcounter{utterance}  
    & & \multicolumn{4}{p{0.6\linewidth}}{
        \cellcolor[rgb]{0.9,0.9,0.9}{
            \makecell[{{p{\linewidth}}}]{
                \texttt{\tiny{[GM$\rangle$P2]}}
                \texttt{Mesajul dvs. a fost transmis celuilalt jucător.} \\
\\ 
\\ 
\texttt{În prezent jucați runda 3 din maximum 12.} \\
\texttt{Ați acumulat în mod colectiv 1 din 8 penalizări.} \\
\texttt{Celălalt jucător a trimis acest mesaj:} \\
\\ 
\\ 
\texttt{"(4,3) e zid pe tabla mea. Încerc (4,4) pentru L."} \\
\\ 
\\ 
\texttt{Care este următoarea dvs. comandă?} \\
            }
        }
    }
    & \\ \\

    \theutterance \stepcounter{utterance}  
    & & & \multicolumn{4}{p{0.6\linewidth}}{
        \cellcolor[rgb]{0.9,0.9,0.9}{
            \makecell[{{p{\linewidth}}}]{
                \texttt{\tiny{[GM$\langle$P2]}}
                \texttt{SPUNE: Înțeles. Pe tabla mea (4,3) e liber. Dacă pentru tine (4,4) e ok pentru L, akkor mutăm L la (4,4) pe ambele. Eu mut P la (4,4)? Nu, e ocupat. Mut P la (4,5)? Verific... Pe tabla mea (4,5) e liber (doar C e la (6,5)). Dar C trebuie să ajungă la (4,5). Deci P trebuie să fie altundeva. Propun: L{-}>(4,4), C{-}>(4,5), P{-}>(4,3) pe ambele table. Verifică dacă (4,3) și (4,5) sunt libere pe a ta.} \\
            }
        }
    }
     \\ \\

    \theutterance \stepcounter{utterance}  
    & & & \multicolumn{2}{p{0.3\linewidth}}{
        \cellcolor[rgb]{0.9,0.9,0.9}{
            \makecell[{{p{\linewidth}}}]{
                \texttt{\tiny{[GM$|$GM]}}
                \texttt{SPUNE: Înțeles. Pe tabla mea (4,3) e liber. Dacă pentru tine (4,4) e ok pentru L, akkor mutăm L la (4,4) pe ambele. Eu mut P la (4,4)? Nu, e ocupat. Mut P la (4,5)? Verific... Pe tabla mea (4,5) e liber (doar C e la (6,5)). Dar C trebuie să ajungă la (4,5). Deci P trebuie să fie altundeva. Propun: L{-}>(4,4), C{-}>(4,5), P{-}>(4,3) pe ambele table. Verifică dacă (4,3) și (4,5) sunt libere pe a ta.} \\
            }
        }
    }
    & & \\ \\

    \theutterance \stepcounter{utterance}  
    & & \multicolumn{4}{p{0.6\linewidth}}{
        \cellcolor[rgb]{0.9,0.9,0.9}{
            \makecell[{{p{\linewidth}}}]{
                \texttt{\tiny{[P1$\langle$GM]}}
                \texttt{Mesajul dvs. a fost transmis celuilalt jucător.} \\
\\ 
\\ 
\texttt{În prezent jucați runda 3 din maximum 12.} \\
\texttt{Ați acumulat în mod colectiv 1 din 8 penalizări.} \\
\texttt{Celălalt jucător a trimis acest mesaj:} \\
\\ 
\\ 
\texttt{"Înțeles. Pe tabla mea (4,3) e liber. Dacă pentru tine (4,4) e ok pentru L, akkor mutăm L la (4,4) pe ambele. Eu mut P la (4,4)? Nu, e ocupat. Mut P la (4,5)? Verific... Pe tabla mea (4,5) e liber (doar C e la (6,5)). Dar C trebuie să ajungă la (4,5). Deci P trebuie să fie altundeva. Propun: L{-}>(4,4), C{-}>(4,5), P{-}>(4,3) pe ambele table. Verifică dacă (4,3) și (4,5) sunt libere pe a ta."} \\
\\ 
\\ 
\texttt{Care este următoarea dvs. comandă?} \\
            }
        }
    }
    & \\ \\

    \theutterance \stepcounter{utterance}  
    & \multicolumn{4}{p{0.6\linewidth}}{
        \cellcolor[rgb]{0.9,0.9,0.9}{
            \makecell[{{p{\linewidth}}}]{
                \texttt{\tiny{[P1$\rangle$GM]}}
                \texttt{MUTĂ: L, (4, 4)} \\
            }
        }
    }
    & & \\ \\

    \theutterance \stepcounter{utterance}  
    & & & \multicolumn{2}{p{0.3\linewidth}}{
        \cellcolor[rgb]{0.9,0.9,0.9}{
            \makecell[{{p{\linewidth}}}]{
                \texttt{\tiny{[GM$|$GM]}}
                \texttt{MUTĂ: L, (4, 4)} \\
            }
        }
    }
    & & \\ \\

    \theutterance \stepcounter{utterance}  
    & & & \multicolumn{2}{p{0.3\linewidth}}{
        \cellcolor[rgb]{0.9,0.9,0.9}{
            \makecell[{{p{\linewidth}}}]{
                \texttt{\tiny{[GM$|$GM]}}
                \texttt{Penalizare: (4,4) nu este gol, ci conține '│'.} \\
\\ 
\texttt{Ați acumulat în mod colectiv 2 din 8 penalizări.} \\
\texttt{Vă rugăm să încercați din nou!} \\
            }
        }
    }
    & & \\ \\

    \theutterance \stepcounter{utterance}  
    & & \multicolumn{4}{p{0.6\linewidth}}{
        \cellcolor[rgb]{0.9,0.9,0.9}{
            \makecell[{{p{\linewidth}}}]{
                \texttt{\tiny{[P1$\langle$GM]}}
                \texttt{Penalizare: (4,4) nu este gol, ci conține '│'.} \\
\\ 
\texttt{Ați acumulat în mod colectiv 2 din 8 penalizări.} \\
\texttt{Vă rugăm să încercați din nou!} \\
            }
        }
    }
    & \\ \\

    \theutterance \stepcounter{utterance}  
    & \multicolumn{4}{p{0.6\linewidth}}{
        \cellcolor[rgb]{0.9,0.9,0.9}{
            \makecell[{{p{\linewidth}}}]{
                \texttt{\tiny{[P1$\rangle$GM]}}
                \texttt{SPUNE: (4,4) e tot zid. Pe coloana 4 doar (4,6) e liber pe tabla mea. Coloana 3 are (3,3) liber. Propun: L la (3,3), C la (3,5), P la (3,4). Verifică dacă sunt libere pe a ta.} \\
            }
        }
    }
    & & \\ \\

    \theutterance \stepcounter{utterance}  
    & & & \multicolumn{2}{p{0.3\linewidth}}{
        \cellcolor[rgb]{0.9,0.9,0.9}{
            \makecell[{{p{\linewidth}}}]{
                \texttt{\tiny{[GM$|$GM]}}
                \texttt{SPUNE: (4,4) e tot zid. Pe coloana 4 doar (4,6) e liber pe tabla mea. Coloana 3 are (3,3) liber. Propun: L la (3,3), C la (3,5), P la (3,4). Verifică dacă sunt libere pe a ta.} \\
            }
        }
    }
    & & \\ \\

    \theutterance \stepcounter{utterance}  
    & & \multicolumn{4}{p{0.6\linewidth}}{
        \cellcolor[rgb]{0.9,0.9,0.9}{
            \makecell[{{p{\linewidth}}}]{
                \texttt{\tiny{[GM$\rangle$P2]}}
                \texttt{Mesajul dvs. a fost transmis celuilalt jucător.} \\
\\ 
\\ 
\texttt{În prezent jucați runda 4 din maximum 12.} \\
\texttt{Ați acumulat în mod colectiv 2 din 8 penalizări.} \\
\texttt{Celălalt jucător a trimis acest mesaj:} \\
\\ 
\\ 
\texttt{"(4,4) e tot zid. Pe coloana 4 doar (4,6) e liber pe tabla mea. Coloana 3 are (3,3) liber. Propun: L la (3,3), C la (3,5), P la (3,4). Verifică dacă sunt libere pe a ta."} \\
\\ 
\\ 
\texttt{Care este următoarea dvs. comandă?} \\
            }
        }
    }
    & \\ \\

    \theutterance \stepcounter{utterance}  
    & & & \multicolumn{4}{p{0.6\linewidth}}{
        \cellcolor[rgb]{0.9,0.9,0.9}{
            \makecell[{{p{\linewidth}}}]{
                \texttt{\tiny{[GM$\langle$P2]}}
                \texttt{SPUNE: Verificat pe tabla mea: (3,3), (3,4), (3,5) sunt toate libere. Sunt de acord cu noua configurație: L la (3,3), C la (3,5), P la (3,4). Încep mutările. Vă anunț după fiecare pas. MUTĂ: L, (3,3)} \\
            }
        }
    }
     \\ \\

    \theutterance \stepcounter{utterance}  
    & & & \multicolumn{2}{p{0.3\linewidth}}{
        \cellcolor[rgb]{0.9,0.9,0.9}{
            \makecell[{{p{\linewidth}}}]{
                \texttt{\tiny{[GM$|$GM]}}
                \texttt{SPUNE: Verificat pe tabla mea: (3,3), (3,4), (3,5) sunt toate libere. Sunt de acord cu noua configurație: L la (3,3), C la (3,5), P la (3,4). Încep mutările. Vă anunț după fiecare pas. MUTĂ: L, (3,3)} \\
            }
        }
    }
    & & \\ \\

    \theutterance \stepcounter{utterance}  
    & & & \multicolumn{2}{p{0.3\linewidth}}{
        \cellcolor[rgb]{0.9,0.9,0.9}{
            \makecell[{{p{\linewidth}}}]{
                \texttt{\tiny{[GM$|$GM]}}
                \texttt{Invalid response format: SPUNE: Verificat pe tabla mea: (3,3), (3,4), (3,5) sunt toate libere. Sunt de acord cu noua configurație: L la (3,3), C la (3,5), P la (3,4). Încep mutările. Vă anunț după fiecare pas. MUTĂ: L, (3,3)} \\
            }
        }
    }
    & & \\ \\

    \theutterance \stepcounter{utterance}  
    & & \multicolumn{4}{p{0.6\linewidth}}{
        \cellcolor[rgb]{0.9,0.9,0.9}{
            \makecell[{{p{\linewidth}}}]{
                \texttt{\tiny{[GM$\rangle$P2]}}
                \texttt{Penalizare: Mesajul dvs. conține mai mult de o comandă!} \\
\\ 
\texttt{Asigurați{-}vă că răspunsul dvs. conține doar `SPUNE: <MESAJ>` sau `MUTĂ: <OBIECT>, (<X>, <Y>)`, și nimic altceva!} \\
\\ 
\texttt{Ați acumulat în mod colectiv 3 din 8 penalizări.} \\
\texttt{Vă rugăm să încercați din nou!} \\
            }
        }
    }
    & \\ \\

    \theutterance \stepcounter{utterance}  
    & & & \multicolumn{4}{p{0.6\linewidth}}{
        \cellcolor[rgb]{0.9,0.9,0.9}{
            \makecell[{{p{\linewidth}}}]{
                \texttt{\tiny{[GM$\langle$P2]}}
                \texttt{MUTĂ: L, (3,3)} \\
            }
        }
    }
     \\ \\

    \theutterance \stepcounter{utterance}  
    & & & \multicolumn{2}{p{0.3\linewidth}}{
        \cellcolor[rgb]{0.9,0.9,0.9}{
            \makecell[{{p{\linewidth}}}]{
                \texttt{\tiny{[GM$|$GM]}}
                \texttt{MUTĂ: L, (3,3)} \\
            }
        }
    }
    & & \\ \\

    \theutterance \stepcounter{utterance}  
    & & & \multicolumn{2}{p{0.3\linewidth}}{
        \cellcolor[rgb]{0.9,0.9,0.9}{
            \makecell[{{p{\linewidth}}}]{
                \texttt{\tiny{[GM$|$GM]}}
                \texttt{Penalizare: (3,3) nu este gol, ci conține 'P'.} \\
\\ 
\texttt{Ați acumulat în mod colectiv 4 din 8 penalizări.} \\
\texttt{Vă rugăm să încercați din nou!} \\
            }
        }
    }
    & & \\ \\

    \theutterance \stepcounter{utterance}  
    & & \multicolumn{4}{p{0.6\linewidth}}{
        \cellcolor[rgb]{0.9,0.9,0.9}{
            \makecell[{{p{\linewidth}}}]{
                \texttt{\tiny{[GM$\rangle$P2]}}
                \texttt{Penalizare: (3,3) nu este gol, ci conține 'P'.} \\
\\ 
\texttt{Ați acumulat în mod colectiv 4 din 8 penalizări.} \\
\texttt{Vă rugăm să încercați din nou!} \\
            }
        }
    }
    & \\ \\

    \theutterance \stepcounter{utterance}  
    & & & \multicolumn{4}{p{0.6\linewidth}}{
        \cellcolor[rgb]{0.9,0.9,0.9}{
            \makecell[{{p{\linewidth}}}]{
                \texttt{\tiny{[GM$\langle$P2]}}
                \texttt{SPUNE: Eroare la mutarea L: (3,3) e ocupat de P pe tabla mea. Mut mai întâi P{-}ul. MUTĂ: P, (3,4)} \\
            }
        }
    }
     \\ \\

    \theutterance \stepcounter{utterance}  
    & & & \multicolumn{2}{p{0.3\linewidth}}{
        \cellcolor[rgb]{0.9,0.9,0.9}{
            \makecell[{{p{\linewidth}}}]{
                \texttt{\tiny{[GM$|$GM]}}
                \texttt{SPUNE: Eroare la mutarea L: (3,3) e ocupat de P pe tabla mea. Mut mai întâi P{-}ul. MUTĂ: P, (3,4)} \\
            }
        }
    }
    & & \\ \\

    \theutterance \stepcounter{utterance}  
    & & & \multicolumn{2}{p{0.3\linewidth}}{
        \cellcolor[rgb]{0.9,0.9,0.9}{
            \makecell[{{p{\linewidth}}}]{
                \texttt{\tiny{[GM$|$GM]}}
                \texttt{Invalid response format: SPUNE: Eroare la mutarea L: (3,3) e ocupat de P pe tabla mea. Mut mai întâi P{-}ul. MUTĂ: P, (3,4)} \\
            }
        }
    }
    & & \\ \\

    \theutterance \stepcounter{utterance}  
    & & \multicolumn{4}{p{0.6\linewidth}}{
        \cellcolor[rgb]{0.9,0.9,0.9}{
            \makecell[{{p{\linewidth}}}]{
                \texttt{\tiny{[GM$\rangle$P2]}}
                \texttt{Penalizare: Mesajul dvs. conține mai mult de o comandă!} \\
\\ 
\texttt{Asigurați{-}vă că răspunsul dvs. conține doar `SPUNE: <MESAJ>` sau `MUTĂ: <OBIECT>, (<X>, <Y>)`, și nimic altceva!} \\
\\ 
\texttt{Ați acumulat în mod colectiv 5 din 8 penalizări.} \\
\texttt{Vă rugăm să încercați din nou!} \\
            }
        }
    }
    & \\ \\

    \theutterance \stepcounter{utterance}  
    & & & \multicolumn{4}{p{0.6\linewidth}}{
        \cellcolor[rgb]{0.9,0.9,0.9}{
            \makecell[{{p{\linewidth}}}]{
                \texttt{\tiny{[GM$\langle$P2]}}
                \texttt{MUTĂ: P, (3,4)} \\
            }
        }
    }
     \\ \\

    \theutterance \stepcounter{utterance}  
    & & & \multicolumn{2}{p{0.3\linewidth}}{
        \cellcolor[rgb]{0.9,0.9,0.9}{
            \makecell[{{p{\linewidth}}}]{
                \texttt{\tiny{[GM$|$GM]}}
                \texttt{MUTĂ: P, (3,4)} \\
            }
        }
    }
    & & \\ \\

    \theutterance \stepcounter{utterance}  
    & & & \multicolumn{2}{p{0.3\linewidth}}{
        \cellcolor[rgb]{0.9,0.9,0.9}{
            \makecell[{{p{\linewidth}}}]{
                \texttt{\tiny{[GM$|$GM]}}
                \texttt{P moved to (3, 4)} \\
            }
        }
    }
    & & \\ \\

    \theutterance \stepcounter{utterance}  
    & & \multicolumn{4}{p{0.6\linewidth}}{
        \cellcolor[rgb]{0.9,0.9,0.9}{
            \makecell[{{p{\linewidth}}}]{
                \texttt{\tiny{[P1$\langle$GM]}}
                \texttt{Mesajul dvs. a fost transmis celuilalt jucător.} \\
\\ 
\\ 
\texttt{În prezent jucați runda 4 din maximum 12.} \\
\texttt{Ați acumulat în mod colectiv 5 din 8 penalizări.} \\
\texttt{Celălalt jucător a mutat un obiect pe grila sa.} \\
\\ 
\\ 
\texttt{Care este următoarea dvs. comandă?} \\
            }
        }
    }
    & \\ \\

    \theutterance \stepcounter{utterance}  
    & \multicolumn{4}{p{0.6\linewidth}}{
        \cellcolor[rgb]{0.9,0.9,0.9}{
            \makecell[{{p{\linewidth}}}]{
                \texttt{\tiny{[P1$\rangle$GM]}}
                \texttt{SPUNE: Ai reușit să muți ceva? Pe tabla mea (3,3) e liber. încerc să mut L acolo.} \\
            }
        }
    }
    & & \\ \\

    \theutterance \stepcounter{utterance}  
    & & & \multicolumn{2}{p{0.3\linewidth}}{
        \cellcolor[rgb]{0.9,0.9,0.9}{
            \makecell[{{p{\linewidth}}}]{
                \texttt{\tiny{[GM$|$GM]}}
                \texttt{SPUNE: Ai reușit să muți ceva? Pe tabla mea (3,3) e liber. încerc să mut L acolo.} \\
            }
        }
    }
    & & \\ \\

    \theutterance \stepcounter{utterance}  
    & & \multicolumn{4}{p{0.6\linewidth}}{
        \cellcolor[rgb]{0.9,0.9,0.9}{
            \makecell[{{p{\linewidth}}}]{
                \texttt{\tiny{[GM$\rangle$P2]}}
                \texttt{Obiectul 'P' a fost mutat la (3,4) cu succes. Grila dvs. actualizată arată astfel:} \\
\texttt{```} \\
\texttt{ 1234567} \\
\texttt{╔═╤═╤═══╗} \\
\texttt{║◌│◌│◌◌◌║ 1} \\
\texttt{║◌├─┤◌L◌║ 2} \\
\texttt{╟─┤◌├─┬─╢ 3} \\
\texttt{║◌│P│◌│◌║ 4} \\
\texttt{╟─┼─┘C│◌║ 5} \\
\texttt{║◌│◌◌┌┴─╢ 6} \\
\texttt{║◌│◌◌│◌◌║ 7} \\
\texttt{╚═╧══╧══╝} \\
\texttt{```} \\
\texttt{În prezent jucați runda 5 din maximum 12.} \\
\texttt{Ați acumulat în mod colectiv 5 din 8 penalizări.} \\
\texttt{Celălalt jucător a trimis acest mesaj:} \\
\\ 
\\ 
\texttt{"Ai reușit să muți ceva? Pe tabla mea (3,3) e liber. încerc să mut L acolo."} \\
\\ 
\\ 
\texttt{Care este următoarea dvs. comandă?} \\
            }
        }
    }
    & \\ \\

    \theutterance \stepcounter{utterance}  
    & & & \multicolumn{4}{p{0.6\linewidth}}{
        \cellcolor[rgb]{0.9,0.9,0.9}{
            \makecell[{{p{\linewidth}}}]{
                \texttt{\tiny{[GM$\langle$P2]}}
                \texttt{MUTĂ: L, (3,3)} \\
            }
        }
    }
     \\ \\

    \theutterance \stepcounter{utterance}  
    & & & \multicolumn{2}{p{0.3\linewidth}}{
        \cellcolor[rgb]{0.9,0.9,0.9}{
            \makecell[{{p{\linewidth}}}]{
                \texttt{\tiny{[GM$|$GM]}}
                \texttt{MUTĂ: L, (3,3)} \\
            }
        }
    }
    & & \\ \\

    \theutterance \stepcounter{utterance}  
    & & & \multicolumn{2}{p{0.3\linewidth}}{
        \cellcolor[rgb]{0.9,0.9,0.9}{
            \makecell[{{p{\linewidth}}}]{
                \texttt{\tiny{[GM$|$GM]}}
                \texttt{L moved to (3, 3)} \\
            }
        }
    }
    & & \\ \\

    \theutterance \stepcounter{utterance}  
    & & \multicolumn{4}{p{0.6\linewidth}}{
        \cellcolor[rgb]{0.9,0.9,0.9}{
            \makecell[{{p{\linewidth}}}]{
                \texttt{\tiny{[P1$\langle$GM]}}
                \texttt{Mesajul dvs. a fost transmis celuilalt jucător.} \\
\\ 
\\ 
\texttt{În prezent jucați runda 5 din maximum 12.} \\
\texttt{Ați acumulat în mod colectiv 5 din 8 penalizări.} \\
\texttt{Celălalt jucător a mutat un obiect pe grila sa.} \\
\\ 
\\ 
\texttt{Care este următoarea dvs. comandă?} \\
            }
        }
    }
    & \\ \\

    \theutterance \stepcounter{utterance}  
    & \multicolumn{4}{p{0.6\linewidth}}{
        \cellcolor[rgb]{0.9,0.9,0.9}{
            \makecell[{{p{\linewidth}}}]{
                \texttt{\tiny{[P1$\rangle$GM]}}
                \texttt{MUTĂ: L, (3, 3)} \\
            }
        }
    }
    & & \\ \\

    \theutterance \stepcounter{utterance}  
    & & & \multicolumn{2}{p{0.3\linewidth}}{
        \cellcolor[rgb]{0.9,0.9,0.9}{
            \makecell[{{p{\linewidth}}}]{
                \texttt{\tiny{[GM$|$GM]}}
                \texttt{MUTĂ: L, (3, 3)} \\
            }
        }
    }
    & & \\ \\

    \theutterance \stepcounter{utterance}  
    & & & \multicolumn{2}{p{0.3\linewidth}}{
        \cellcolor[rgb]{0.9,0.9,0.9}{
            \makecell[{{p{\linewidth}}}]{
                \texttt{\tiny{[GM$|$GM]}}
                \texttt{L moved to (3, 3)} \\
            }
        }
    }
    & & \\ \\

    \theutterance \stepcounter{utterance}  
    & & \multicolumn{4}{p{0.6\linewidth}}{
        \cellcolor[rgb]{0.9,0.9,0.9}{
            \makecell[{{p{\linewidth}}}]{
                \texttt{\tiny{[GM$\rangle$P2]}}
                \texttt{Obiectul 'L' a fost mutat la (3,3) cu succes. Grila dvs. actualizată arată astfel:} \\
\texttt{```} \\
\texttt{ 1234567} \\
\texttt{╔═╤═╤═══╗} \\
\texttt{║◌│◌│◌◌◌║ 1} \\
\texttt{║◌├─┤◌◌◌║ 2} \\
\texttt{╟─┤L├─┬─╢ 3} \\
\texttt{║◌│P│◌│◌║ 4} \\
\texttt{╟─┼─┘C│◌║ 5} \\
\texttt{║◌│◌◌┌┴─╢ 6} \\
\texttt{║◌│◌◌│◌◌║ 7} \\
\texttt{╚═╧══╧══╝} \\
\texttt{```} \\
\texttt{În prezent jucați runda 6 din maximum 12.} \\
\texttt{Ați acumulat în mod colectiv 5 din 8 penalizări.} \\
\texttt{Celălalt jucător a mutat un obiect pe grila sa.} \\
\\ 
\\ 
\texttt{Care este următoarea dvs. comandă?} \\
            }
        }
    }
    & \\ \\

    \theutterance \stepcounter{utterance}  
    & & & \multicolumn{4}{p{0.6\linewidth}}{
        \cellcolor[rgb]{0.9,0.9,0.9}{
            \makecell[{{p{\linewidth}}}]{
                \texttt{\tiny{[GM$\langle$P2]}}
                \texttt{MUTĂ: C, (3,5)} \\
            }
        }
    }
     \\ \\

    \theutterance \stepcounter{utterance}  
    & & & \multicolumn{2}{p{0.3\linewidth}}{
        \cellcolor[rgb]{0.9,0.9,0.9}{
            \makecell[{{p{\linewidth}}}]{
                \texttt{\tiny{[GM$|$GM]}}
                \texttt{MUTĂ: C, (3,5)} \\
            }
        }
    }
    & & \\ \\

    \theutterance \stepcounter{utterance}  
    & & & \multicolumn{2}{p{0.3\linewidth}}{
        \cellcolor[rgb]{0.9,0.9,0.9}{
            \makecell[{{p{\linewidth}}}]{
                \texttt{\tiny{[GM$|$GM]}}
                \texttt{Penalizare: (3,5) nu este gol, ci conține '─'.} \\
\\ 
\texttt{Ați acumulat în mod colectiv 6 din 8 penalizări.} \\
\texttt{Vă rugăm să încercați din nou!} \\
            }
        }
    }
    & & \\ \\

    \theutterance \stepcounter{utterance}  
    & & \multicolumn{4}{p{0.6\linewidth}}{
        \cellcolor[rgb]{0.9,0.9,0.9}{
            \makecell[{{p{\linewidth}}}]{
                \texttt{\tiny{[GM$\rangle$P2]}}
                \texttt{Penalizare: (3,5) nu este gol, ci conține '─'.} \\
\\ 
\texttt{Ați acumulat în mod colectiv 6 din 8 penalizări.} \\
\texttt{Vă rugăm să încercați din nou!} \\
            }
        }
    }
    & \\ \\

    \theutterance \stepcounter{utterance}  
    & & & \multicolumn{4}{p{0.6\linewidth}}{
        \cellcolor[rgb]{0.9,0.9,0.9}{
            \makecell[{{p{\linewidth}}}]{
                \texttt{\tiny{[GM$\langle$P2]}}
                \texttt{SPUNE: Eroare: (3,5) e perete pe tabla mea. C e la (6,5). Propun C la (3,6) dacă e liber pe a ta. Verifică (3,6) te rog.} \\
            }
        }
    }
     \\ \\

    \theutterance \stepcounter{utterance}  
    & & & \multicolumn{2}{p{0.3\linewidth}}{
        \cellcolor[rgb]{0.9,0.9,0.9}{
            \makecell[{{p{\linewidth}}}]{
                \texttt{\tiny{[GM$|$GM]}}
                \texttt{SPUNE: Eroare: (3,5) e perete pe tabla mea. C e la (6,5). Propun C la (3,6) dacă e liber pe a ta. Verifică (3,6) te rog.} \\
            }
        }
    }
    & & \\ \\

    \theutterance \stepcounter{utterance}  
    & & \multicolumn{4}{p{0.6\linewidth}}{
        \cellcolor[rgb]{0.9,0.9,0.9}{
            \makecell[{{p{\linewidth}}}]{
                \texttt{\tiny{[P1$\langle$GM]}}
                \texttt{Obiectul 'L' a fost mutat la (3,3) cu succes. Grila dvs. actualizată arată astfel:} \\
\texttt{```} \\
\texttt{ 1234567} \\
\texttt{╔═╤═╤═══╗} \\
\texttt{║◌│◌│◌◌◌║ 1} \\
\texttt{║C├─┤◌◌◌║ 2} \\
\texttt{╟─┤L├─┬─╢ 3} \\
\texttt{║◌│◌│◌│◌║ 4} \\
\texttt{╟─┼─┘◌│◌║ 5} \\
\texttt{║◌│◌◌┌┴─╢ 6} \\
\texttt{║◌│◌◌│◌P║ 7} \\
\texttt{╚═╧══╧══╝} \\
\texttt{```} \\
\texttt{În prezent jucați runda 6 din maximum 12.} \\
\texttt{Ați acumulat în mod colectiv 6 din 8 penalizări.} \\
\texttt{Celălalt jucător a trimis acest mesaj:} \\
\\ 
\\ 
\texttt{"Eroare: (3,5) e perete pe tabla mea. C e la (6,5). Propun C la (3,6) dacă e liber pe a ta. Verifică (3,6) te rog."} \\
\\ 
\\ 
\texttt{Care este următoarea dvs. comandă?} \\
            }
        }
    }
    & \\ \\

    \theutterance \stepcounter{utterance}  
    & \multicolumn{4}{p{0.6\linewidth}}{
        \cellcolor[rgb]{0.9,0.9,0.9}{
            \makecell[{{p{\linewidth}}}]{
                \texttt{\tiny{[P1$\rangle$GM]}}
                \texttt{SPUNE: (3,6) e liber pe a mea. Mut C la (3,6) pe ambele. Tu mută C{-}ul tău de la (6,5) la (3,6). Eu mut C{-}ul de la (1,2) la (3,6).} \\
            }
        }
    }
    & & \\ \\

    \theutterance \stepcounter{utterance}  
    & & & \multicolumn{2}{p{0.3\linewidth}}{
        \cellcolor[rgb]{0.9,0.9,0.9}{
            \makecell[{{p{\linewidth}}}]{
                \texttt{\tiny{[GM$|$GM]}}
                \texttt{SPUNE: (3,6) e liber pe a mea. Mut C la (3,6) pe ambele. Tu mută C{-}ul tău de la (6,5) la (3,6). Eu mut C{-}ul de la (1,2) la (3,6).} \\
            }
        }
    }
    & & \\ \\

    \theutterance \stepcounter{utterance}  
    & & \multicolumn{4}{p{0.6\linewidth}}{
        \cellcolor[rgb]{0.9,0.9,0.9}{
            \makecell[{{p{\linewidth}}}]{
                \texttt{\tiny{[GM$\rangle$P2]}}
                \texttt{Mesajul dvs. a fost transmis celuilalt jucător.} \\
\\ 
\\ 
\texttt{În prezent jucați runda 7 din maximum 12.} \\
\texttt{Ați acumulat în mod colectiv 6 din 8 penalizări.} \\
\texttt{Celălalt jucător a trimis acest mesaj:} \\
\\ 
\\ 
\texttt{"(3,6) e liber pe a mea. Mut C la (3,6) pe ambele. Tu mută C{-}ul tău de la (6,5) la (3,6). Eu mut C{-}ul de la (1,2) la (3,6)."} \\
\\ 
\\ 
\texttt{Care este următoarea dvs. comandă?} \\
            }
        }
    }
    & \\ \\

    \theutterance \stepcounter{utterance}  
    & & & \multicolumn{4}{p{0.6\linewidth}}{
        \cellcolor[rgb]{0.9,0.9,0.9}{
            \makecell[{{p{\linewidth}}}]{
                \texttt{\tiny{[GM$\langle$P2]}}
                \texttt{MUTĂ: C, (3,6)} \\
            }
        }
    }
     \\ \\

    \theutterance \stepcounter{utterance}  
    & & & \multicolumn{2}{p{0.3\linewidth}}{
        \cellcolor[rgb]{0.9,0.9,0.9}{
            \makecell[{{p{\linewidth}}}]{
                \texttt{\tiny{[GM$|$GM]}}
                \texttt{MUTĂ: C, (3,6)} \\
            }
        }
    }
    & & \\ \\

    \theutterance \stepcounter{utterance}  
    & & & \multicolumn{2}{p{0.3\linewidth}}{
        \cellcolor[rgb]{0.9,0.9,0.9}{
            \makecell[{{p{\linewidth}}}]{
                \texttt{\tiny{[GM$|$GM]}}
                \texttt{C moved to (3, 6)} \\
            }
        }
    }
    & & \\ \\

    \theutterance \stepcounter{utterance}  
    & & \multicolumn{4}{p{0.6\linewidth}}{
        \cellcolor[rgb]{0.9,0.9,0.9}{
            \makecell[{{p{\linewidth}}}]{
                \texttt{\tiny{[P1$\langle$GM]}}
                \texttt{Mesajul dvs. a fost transmis celuilalt jucător.} \\
\\ 
\\ 
\texttt{În prezent jucați runda 7 din maximum 12.} \\
\texttt{Ați acumulat în mod colectiv 6 din 8 penalizări.} \\
\texttt{Celălalt jucător a mutat un obiect pe grila sa.} \\
\\ 
\\ 
\texttt{Care este următoarea dvs. comandă?} \\
            }
        }
    }
    & \\ \\

    \theutterance \stepcounter{utterance}  
    & \multicolumn{4}{p{0.6\linewidth}}{
        \cellcolor[rgb]{0.9,0.9,0.9}{
            \makecell[{{p{\linewidth}}}]{
                \texttt{\tiny{[P1$\rangle$GM]}}
                \texttt{MUTĂ: C, (3, 6)} \\
            }
        }
    }
    & & \\ \\

    \theutterance \stepcounter{utterance}  
    & & & \multicolumn{2}{p{0.3\linewidth}}{
        \cellcolor[rgb]{0.9,0.9,0.9}{
            \makecell[{{p{\linewidth}}}]{
                \texttt{\tiny{[GM$|$GM]}}
                \texttt{MUTĂ: C, (3, 6)} \\
            }
        }
    }
    & & \\ \\

    \theutterance \stepcounter{utterance}  
    & & & \multicolumn{2}{p{0.3\linewidth}}{
        \cellcolor[rgb]{0.9,0.9,0.9}{
            \makecell[{{p{\linewidth}}}]{
                \texttt{\tiny{[GM$|$GM]}}
                \texttt{C moved to (3, 6)} \\
            }
        }
    }
    & & \\ \\

    \theutterance \stepcounter{utterance}  
    & & \multicolumn{4}{p{0.6\linewidth}}{
        \cellcolor[rgb]{0.9,0.9,0.9}{
            \makecell[{{p{\linewidth}}}]{
                \texttt{\tiny{[GM$\rangle$P2]}}
                \texttt{Obiectul 'C' a fost mutat la (3,6) cu succes. Grila dvs. actualizată arată astfel:} \\
\texttt{```} \\
\texttt{ 1234567} \\
\texttt{╔═╤═╤═══╗} \\
\texttt{║◌│◌│◌◌◌║ 1} \\
\texttt{║◌├─┤◌◌◌║ 2} \\
\texttt{╟─┤L├─┬─╢ 3} \\
\texttt{║◌│P│◌│◌║ 4} \\
\texttt{╟─┼─┘◌│◌║ 5} \\
\texttt{║◌│C◌┌┴─╢ 6} \\
\texttt{║◌│◌◌│◌◌║ 7} \\
\texttt{╚═╧══╧══╝} \\
\texttt{```} \\
\texttt{În prezent jucați runda 8 din maximum 12.} \\
\texttt{Ați acumulat în mod colectiv 6 din 8 penalizări.} \\
\texttt{Celălalt jucător a mutat un obiect pe grila sa.} \\
\\ 
\\ 
\texttt{Care este următoarea dvs. comandă?} \\
            }
        }
    }
    & \\ \\

    \theutterance \stepcounter{utterance}  
    & & & \multicolumn{4}{p{0.6\linewidth}}{
        \cellcolor[rgb]{0.9,0.9,0.9}{
            \makecell[{{p{\linewidth}}}]{
                \texttt{\tiny{[GM$\langle$P2]}}
                \texttt{SPUNE: Toate obiectele mele sunt în poziție finală: L(3,3), P(3,4), C(3,6). Confirmă când ai terminat și tu.} \\
            }
        }
    }
     \\ \\

    \theutterance \stepcounter{utterance}  
    & & & \multicolumn{2}{p{0.3\linewidth}}{
        \cellcolor[rgb]{0.9,0.9,0.9}{
            \makecell[{{p{\linewidth}}}]{
                \texttt{\tiny{[GM$|$GM]}}
                \texttt{SPUNE: Toate obiectele mele sunt în poziție finală: L(3,3), P(3,4), C(3,6). Confirmă când ai terminat și tu.} \\
            }
        }
    }
    & & \\ \\

    \theutterance \stepcounter{utterance}  
    & & \multicolumn{4}{p{0.6\linewidth}}{
        \cellcolor[rgb]{0.9,0.9,0.9}{
            \makecell[{{p{\linewidth}}}]{
                \texttt{\tiny{[P1$\langle$GM]}}
                \texttt{Obiectul 'C' a fost mutat la (3,6) cu succes. Grila dvs. actualizată arată astfel:} \\
\texttt{```} \\
\texttt{ 1234567} \\
\texttt{╔═╤═╤═══╗} \\
\texttt{║◌│◌│◌◌◌║ 1} \\
\texttt{║◌├─┤◌◌◌║ 2} \\
\texttt{╟─┤L├─┬─╢ 3} \\
\texttt{║◌│◌│◌│◌║ 4} \\
\texttt{╟─┼─┘◌│◌║ 5} \\
\texttt{║◌│C◌┌┴─╢ 6} \\
\texttt{║◌│◌◌│◌P║ 7} \\
\texttt{╚═╧══╧══╝} \\
\texttt{```} \\
\texttt{În prezent jucați runda 8 din maximum 12.} \\
\texttt{Ați acumulat în mod colectiv 6 din 8 penalizări.} \\
\texttt{Celălalt jucător a trimis acest mesaj:} \\
\\ 
\\ 
\texttt{"Toate obiectele mele sunt în poziție finală: L(3,3), P(3,4), C(3,6). Confirmă când ai terminat și tu."} \\
\\ 
\\ 
\texttt{Care este următoarea dvs. comandă?} \\
            }
        }
    }
    & \\ \\

    \theutterance \stepcounter{utterance}  
    & \multicolumn{4}{p{0.6\linewidth}}{
        \cellcolor[rgb]{0.9,0.9,0.9}{
            \makecell[{{p{\linewidth}}}]{
                \texttt{\tiny{[P1$\rangle$GM]}}
                \texttt{MUTĂ: P, (3, 4)} \\
            }
        }
    }
    & & \\ \\

    \theutterance \stepcounter{utterance}  
    & & & \multicolumn{2}{p{0.3\linewidth}}{
        \cellcolor[rgb]{0.9,0.9,0.9}{
            \makecell[{{p{\linewidth}}}]{
                \texttt{\tiny{[GM$|$GM]}}
                \texttt{MUTĂ: P, (3, 4)} \\
            }
        }
    }
    & & \\ \\

    \theutterance \stepcounter{utterance}  
    & & & \multicolumn{2}{p{0.3\linewidth}}{
        \cellcolor[rgb]{0.9,0.9,0.9}{
            \makecell[{{p{\linewidth}}}]{
                \texttt{\tiny{[GM$|$GM]}}
                \texttt{P moved to (3, 4)} \\
            }
        }
    }
    & & \\ \\

    \theutterance \stepcounter{utterance}  
    & & \multicolumn{4}{p{0.6\linewidth}}{
        \cellcolor[rgb]{0.9,0.9,0.9}{
            \makecell[{{p{\linewidth}}}]{
                \texttt{\tiny{[GM$\rangle$P2]}}
                \texttt{Mesajul dvs. a fost transmis celuilalt jucător.} \\
\\ 
\\ 
\texttt{În prezent jucați runda 9 din maximum 12.} \\
\texttt{Ați acumulat în mod colectiv 6 din 8 penalizări.} \\
\texttt{Celălalt jucător a mutat un obiect pe grila sa.} \\
\\ 
\\ 
\texttt{Care este următoarea dvs. comandă?} \\
            }
        }
    }
    & \\ \\

    \theutterance \stepcounter{utterance}  
    & & & \multicolumn{4}{p{0.6\linewidth}}{
        \cellcolor[rgb]{0.9,0.9,0.9}{
            \makecell[{{p{\linewidth}}}]{
                \texttt{\tiny{[GM$\langle$P2]}}
                \texttt{SPUNE: terminat?} \\
            }
        }
    }
     \\ \\

    \theutterance \stepcounter{utterance}  
    & & & \multicolumn{2}{p{0.3\linewidth}}{
        \cellcolor[rgb]{0.9,0.9,0.9}{
            \makecell[{{p{\linewidth}}}]{
                \texttt{\tiny{[GM$|$GM]}}
                \texttt{SPUNE: terminat?} \\
            }
        }
    }
    & & \\ \\

    \theutterance \stepcounter{utterance}  
    & & \multicolumn{4}{p{0.6\linewidth}}{
        \cellcolor[rgb]{0.9,0.9,0.9}{
            \makecell[{{p{\linewidth}}}]{
                \texttt{\tiny{[P1$\langle$GM]}}
                \texttt{Obiectul 'P' a fost mutat la (3,4) cu succes. Grila dvs. actualizată arată astfel:} \\
\texttt{```} \\
\texttt{ 1234567} \\
\texttt{╔═╤═╤═══╗} \\
\texttt{║◌│◌│◌◌◌║ 1} \\
\texttt{║◌├─┤◌◌◌║ 2} \\
\texttt{╟─┤L├─┬─╢ 3} \\
\texttt{║◌│P│◌│◌║ 4} \\
\texttt{╟─┼─┘◌│◌║ 5} \\
\texttt{║◌│C◌┌┴─╢ 6} \\
\texttt{║◌│◌◌│◌◌║ 7} \\
\texttt{╚═╧══╧══╝} \\
\texttt{```} \\
\texttt{În prezent jucați runda 9 din maximum 12.} \\
\texttt{Ați acumulat în mod colectiv 6 din 8 penalizări.} \\
\texttt{Celălalt jucător a trimis acest mesaj:} \\
\\ 
\\ 
\texttt{"terminat?"} \\
\\ 
\\ 
\texttt{Care este următoarea dvs. comandă?} \\
            }
        }
    }
    & \\ \\

    \theutterance \stepcounter{utterance}  
    & \multicolumn{4}{p{0.6\linewidth}}{
        \cellcolor[rgb]{0.9,0.9,0.9}{
            \makecell[{{p{\linewidth}}}]{
                \texttt{\tiny{[P1$\rangle$GM]}}
                \texttt{SPUNE: terminat!} \\
            }
        }
    }
    & & \\ \\

    \theutterance \stepcounter{utterance}  
    & & & \multicolumn{2}{p{0.3\linewidth}}{
        \cellcolor[rgb]{0.9,0.9,0.9}{
            \makecell[{{p{\linewidth}}}]{
                \texttt{\tiny{[GM$|$GM]}}
                \texttt{SPUNE: terminat!} \\
            }
        }
    }
    & & \\ \\

    \theutterance \stepcounter{utterance}  
    & & & \multicolumn{2}{p{0.3\linewidth}}{
        \cellcolor[rgb]{0.9,0.9,0.9}{
            \makecell[{{p{\linewidth}}}]{
                \texttt{\tiny{[GM$|$GM]}}
                \texttt{true} \\
            }
        }
    }
    & & \\ \\

    \theutterance \stepcounter{utterance}  
    & & & \multicolumn{2}{p{0.3\linewidth}}{
        \cellcolor[rgb]{0.9,0.9,0.9}{
            \makecell[{{p{\linewidth}}}]{
                \texttt{\tiny{[GM$|$GM]}}
                \texttt{Initial states:} \\
\texttt{```} \\
\texttt{Player 1:} \\
\texttt{ 1234567} \\
\texttt{╔═╤═╤═══╗} \\
\texttt{║L│◌│◌◌◌║ 1} \\
\texttt{║C├─┤◌◌◌║ 2} \\
\texttt{╟─┤◌├─┬─╢ 3} \\
\texttt{║◌│◌│◌│◌║ 4} \\
\texttt{╟─┼─┘◌│◌║ 5} \\
\texttt{║◌│◌◌┌┴─╢ 6} \\
\texttt{║◌│◌◌│◌P║ 7} \\
\texttt{╚═╧══╧══╝} \\
\\ 
\\ 
\texttt{Player 2:} \\
\texttt{ 1234567} \\
\texttt{╔═╤═╤═══╗} \\
\texttt{║◌│◌│◌◌◌║ 1} \\
\texttt{║◌├─┤◌L◌║ 2} \\
\texttt{╟─┤P├─┬─╢ 3} \\
\texttt{║◌│◌│◌│◌║ 4} \\
\texttt{╟─┼─┘C│◌║ 5} \\
\texttt{║◌│◌◌┌┴─╢ 6} \\
\texttt{║◌│◌◌│◌◌║ 7} \\
\texttt{╚═╧══╧══╝} \\
\\ 
\texttt{```} \\
            }
        }
    }
    & & \\ \\

    \theutterance \stepcounter{utterance}  
    & & & \multicolumn{2}{p{0.3\linewidth}}{
        \cellcolor[rgb]{0.9,0.9,0.9}{
            \makecell[{{p{\linewidth}}}]{
                \texttt{\tiny{[GM$|$GM]}}
                \texttt{End states:} \\
\texttt{```} \\
\texttt{Player 1:} \\
\texttt{ 1234567} \\
\texttt{╔═╤═╤═══╗} \\
\texttt{║◌│◌│◌◌◌║ 1} \\
\texttt{║◌├─┤◌◌◌║ 2} \\
\texttt{╟─┤L├─┬─╢ 3} \\
\texttt{║◌│P│◌│◌║ 4} \\
\texttt{╟─┼─┘◌│◌║ 5} \\
\texttt{║◌│C◌┌┴─╢ 6} \\
\texttt{║◌│◌◌│◌◌║ 7} \\
\texttt{╚═╧══╧══╝} \\
\\ 
\\ 
\texttt{Player 2:} \\
\texttt{ 1234567} \\
\texttt{╔═╤═╤═══╗} \\
\texttt{║◌│◌│◌◌◌║ 1} \\
\texttt{║◌├─┤◌◌◌║ 2} \\
\texttt{╟─┤L├─┬─╢ 3} \\
\texttt{║◌│P│◌│◌║ 4} \\
\texttt{╟─┼─┘◌│◌║ 5} \\
\texttt{║◌│C◌┌┴─╢ 6} \\
\texttt{║◌│◌◌│◌◌║ 7} \\
\texttt{╚═╧══╧══╝} \\
\\ 
\texttt{```} \\
            }
        }
    }
    & & \\ \\

    \theutterance \stepcounter{utterance}  
    & & & \multicolumn{2}{p{0.3\linewidth}}{
        \cellcolor[rgb]{0.9,0.9,0.9}{
            \makecell[{{p{\linewidth}}}]{
                \texttt{\tiny{[GM$|$GM]}}
                \texttt{* success: True} \\
\texttt{* lose: False} \\
\texttt{* aborted: False} \\
\texttt{{-}{-}{-}{-}{-}{-}{-}} \\
\texttt{* Shifts: 3.00} \\
\texttt{* Max Shifts: 4.00} \\
\texttt{* Min Shifts: 2.00} \\
\texttt{* End Distance Sum: 0.00} \\
\texttt{* Init Distance Sum: 15.76} \\
\texttt{* Expected Distance Sum: 12.57} \\
\texttt{* Penalties: 6.00} \\
\texttt{* Max Penalties: 8.00} \\
\texttt{* Rounds: 9.00} \\
\texttt{* Max Rounds: 12.00} \\
\texttt{* Object Count: 3.00} \\
            }
        }
    }
    & & \\ \\

\end{supertabular}
}

\end{document}
